\documentclass[12pt]{article}
\usepackage[utf8]{inputenc}
\usepackage[T1]{fontenc}
\usepackage{hyperref}
\title{Pwapw\^a Language Resources for K-12 Curriculum Development}
\author{Compiled by Claude Sonnet 5.5 (claude-sonnet-5-5) for Andr\'as Kornai}
\date{}
\begin{document}
\maketitle
\section{Introduction}
Pwapw\^a (ISO 639-3 pop, also Pwapwa, Poapoa) is an Austronesian language of New Caledonia spoken around Voh. It is severely endangered; speaker counts in sources are small and inconsistent (40 in a 2006 census figure cited by Wikipedia, 16 in Verbix). Search found no Pwapw\^a dictionary, textbook or school materials.
\section{Primary Educational Resources}
None found.
\section{Linguistic and Reference Materials}
\begin{itemize}
\item Wikipedia article on Pwapw\^a: \url{https://en.wikipedia.org/wiki/Pwapw%C3%A2_language}
\item Verbix Endangered Languages page for Pwapw\^a: \url{https://docs.verbix.com/EndangeredLanguages/Pwapwa}
\end{itemize}
\section{Religious and Community Resources}
None found specific to Pwapw\^a.
\section{Digital and Audio Resources}
None found. A project of the Acad\'emie des Langues Kanak with Princeton University is producing dictionaries of other Kanak languages (Pa\"ic\^i, C\`emuh\^i, Haeke), not Pwapw\^a: \url{https://outremers360.com/bassin-pacifique-appli/nouvelle-caledonie-des-dictionnaires-en-langues-kanak-bientot-accessibles-au-grand-public-grace-a-un-partenariat-avec-luniversite-de-princeton}
\section{Curriculum Development Context}
Kanak languages in New Caledonia are taught within a French-language school system; the Acad\'emie des Langues Kanak is the natural institutional contact. Any Pwapw\^a curriculum would require new documentation with the remaining speakers.
\end{document}